% ============================================================================
\chapter{2D Physics System}
% ============================================================================

\section{Components}

\subsection{RigidBody2D}

Stores dynamic properties: \texttt{mass} $m$, \texttt{restitution} $e$,
\texttt{friction} $\mu$, \texttt{linearDamping} $d$, flags
\texttt{isKinematic} and \texttt{isSleeping}.

\subsection{Collider2D}

Stores geometric properties: \texttt{shape} (AABB or Circle),
\texttt{size}/\texttt{radius}, \texttt{offset}, collision
\texttt{layer} / \texttt{layerMask} bitmask, and flags
\texttt{isStatic}, \texttt{isTrigger}, \texttt{isOneWay}.

\section{Simulation Loop}

The system runs $k = 2$ sub-steps per frame ($k = \texttt{kSubSteps}$):

\begin{equation}
\Delta t_\text{sub} = \frac{\Delta t}{k}
\end{equation}

Per sub-step:
\begin{enumerate}
  \item Apply queued forces and impulses.
  \item Integrate velocities and positions (semi-implicit Euler).
  \item Build the spatial hash grid.
  \item Detect and resolve collisions.
  \item Update sleep state.
\end{enumerate}

\section{Integration (Semi-Implicit Euler)}

\begin{align}
  \mathbf{v}_{n+1} &= \mathbf{v}_n + \mathbf{g}\,\Delta t_\text{sub}
                    \quad (\text{gravity accumulation}) \\
  \mathbf{v}_{n+1} &\leftarrow \mathbf{v}_{n+1} \cdot \max(0,\, 1 - d\,\Delta t_\text{sub})
                    \quad (\text{linear damping}) \\
  \mathbf{x}_{n+1} &= \mathbf{x}_n + \mathbf{v}_{n+1}\,\Delta t_\text{sub}
\end{align}

The velocity is updated \emph{before} the position (semi-implicit),
which is more stable than explicit Euler for spring-like constraints.

\section{Broad-Phase: Uniform Grid}

The spatial hash maps a 2D integer cell coordinate $(c_x, c_y)$ to a
list of entity IDs. The cell coordinate of a world position $p$ is:

\begin{equation}
c = \left\lfloor \frac{p}{\texttt{kGridCellSize}} \right\rfloor
\end{equation}

with $\texttt{kGridCellSize} = 128$ world units. An entity whose AABB
spans multiple cells is inserted into all cells it touches. The cell key
is computed as:

\begin{equation}
\text{key}(c_x, c_y) = c_x \cdot 73856093 \oplus c_y \cdot 19349663
\end{equation}

(a standard spatial hashing polynomial). Candidate collision pairs are
generated by checking all pairs within the same cell, with duplicate
suppression.

\section{Narrow-Phase Geometry}

\subsection{AABB vs.\ AABB}

Let centres be $\mathbf{p}_A, \mathbf{p}_B$ with half-extents
$\mathbf{h}_A, \mathbf{h}_B$. A collision occurs when:

\begin{align}
  \text{overlapX} &= (h_{Ax} + h_{Bx}) - |p_{Bx} - p_{Ax}| > 0 \\
  \text{overlapY} &= (h_{Ay} + h_{By}) - |p_{By} - p_{Ay}| > 0
\end{align}

The separation axis is the one with \emph{minimum overlap}:

\begin{equation}
\mathbf{n} =
\begin{cases}
  (\text{sgn}(\Delta x),\, 0) & \text{if overlapX} < \text{overlapY} \\
  (0,\, \text{sgn}(\Delta y)) & \text{otherwise}
\end{cases}
\end{equation}

\subsection{Circle vs.\ Circle}

Let distance $d = |\mathbf{p}_B - \mathbf{p}_A|$ and sum of radii
$r = r_A + r_B$.

\begin{align}
  \text{collision} &\Leftrightarrow d < r \\
  \mathbf{n} &= (\mathbf{p}_B - \mathbf{p}_A) / d \\
  \text{penetration} &= r - d
\end{align}

\subsection{Circle vs.\ AABB}

The closest point on the AABB to the circle centre is:

\begin{equation}
\mathbf{c} = \text{clamp}(\mathbf{p}_\text{circle},\,
             \mathbf{p}_\text{AABB} - \mathbf{h},\,
             \mathbf{p}_\text{AABB} + \mathbf{h})
\end{equation}

If $|\mathbf{p}_\text{circle} - \mathbf{c}| < r$, a collision is detected.

\section{Impulse-Based Resolution}

Let relative velocity along the collision normal be:

\begin{equation}
v_\text{rel} = (\mathbf{v}_B - \mathbf{v}_A) \cdot \mathbf{n}
\end{equation}

If $v_\text{rel} > 0$ (separating), no impulse is applied.
Otherwise, the scalar impulse magnitude is:

\begin{equation}
j = \frac{-(1 + e)\, v_\text{rel}}{m_A^{-1} + m_B^{-1}}
\label{eq:impulse}
\end{equation}

where $e = \min(e_A, e_B)$ is the coefficient of restitution.
Velocities are updated:

\begin{align}
  \mathbf{v}_A &\leftarrow \mathbf{v}_A - j\,m_A^{-1}\,\mathbf{n} \\
  \mathbf{v}_B &\leftarrow \mathbf{v}_B + j\,m_B^{-1}\,\mathbf{n}
\end{align}

\subsubsection{Friction}

The tangential impulse $j_t$ is:

\begin{equation}
j_t = \frac{-v_\text{rel,tan}}{m_A^{-1} + m_B^{-1}},
\quad \mu = \sqrt{\mu_A \mu_B}
\end{equation}

Applied as a Coulomb friction cone: $|j_t| \leq \mu |j|$.

\section{Positional Correction (Baumgarte)}

To prevent sinking, the positions are corrected proportionally to the
penetration depth, with a slop tolerance $\delta_\text{slop}$ and
factor $k_\text{B}$:

\begin{equation}
\Delta\mathbf{x} = \frac{(\text{penetration} - \delta_\text{slop}) \cdot k_\text{B}}{m_A^{-1} + m_B^{-1}}\,\mathbf{n}
\end{equation}

with $\delta_\text{slop} = 0.01$ and $k_\text{B} = 0.4$.

\section{Sleep System}

An entity enters sleep when its velocity magnitude drops below
$v_\text{sleep} = 0.05$ for longer than $t_\text{sleep} = 0.5\,\text{s}$.
A sleeping entity skips integration entirely, saving compute for
static scenes with many resting bodies. Any external force or impulse
wakes the entity immediately.

\section{Raycast}

A ray $\mathbf{r}(t) = \mathbf{o} + t\hat{\mathbf{d}}$ is tested against
all entities with a \texttt{Collider2D}. For AABB colliders the
slab method is used; for circle colliders the quadratic form:

\begin{equation}
|(\mathbf{o} + t\hat{\mathbf{d}}) - \mathbf{c}|^2 = r^2
\end{equation}

is solved for $t$. The hit with minimum $t \geq 0$ is returned.
